% Captions generated by run_all.py; all quoted numbers come from output/data/*.csv.
% Mismatch costs are in bits (base-2 logarithms).

% ----------------------------------------------------------------------
% fig04: Ripple-carry adder: gate-by-gate vs layer-by-layer (matched / joint layer priors)
\caption{Comparison of layer-by-layer and gate-by-gate mismatch costs for the
ripple-carry adder circuit family ($n = 2,\dots,5$ bits per operand, circuit size
$|C_n| = 5n$), with uniformly random inputs. Each implementation uses the priors that
minimize its own mismatch cost for this input distribution. For gate-by-gate execution
[Eq.~(26)], every gate is a separate subsystem process with prior
$q_{g|\mathrm{pa}(g)} = p_g$. For layer-by-layer execution [Eq.~(33)], every layer is a
single subsystem process whose joint prior $q_{l|\mathrm{pa}(l)} = p_l$ matches the
correlated distribution of the gates in that layer. The layer-by-layer mismatch cost is
lower by $\sum_l \mathrm{TC}(X_l)$, the total correlation among the gates within each
layer (between 0.76 and 2.46 bits here). If instead each layer's prior is the product
of its gates' priors, as for uniform priors, the two costs coincide.
Mismatch cost in bits.}

% ----------------------------------------------------------------------
% fig06: AND binary tree: all-ones input saturates the upper bound
\caption{Saturating the upper bound on MMC for the AND binary tree circuit family, with
uniform priors for all gates and inputs. The solid blue line is the total mismatch
cost from Eq.~(26) for an input distribution concentrated on the all-ones string. It
coincides with the theoretical upper bound $|C_n| + n$ of Eq.~(C13) (orange dashed):
every gate contributes $\log 2$, and re-initializing the inputs contributes
$D(p_{\mathrm{in}}\|q_{\mathrm{in}}) = n \log 2$. For an i.i.d.\ uniform input
distribution (green), $D(p_{\mathrm{in}}\|q_{\mathrm{in}}) = 0$ and the total mismatch
cost is $|C_n|$ bits, below the bound. With uniform priors, each deterministic gate
contributes exactly one bit for any input distribution, so the two curves differ only
through the input re-initialization term.}

% ----------------------------------------------------------------------
% fig07: Ripple-carry vs Kogge-Stone (type-basis priors and matched priors)
% Option A: one prior per gate type (ripple_vs_kogg2.pdf)
\caption{Mismatch cost of two circuit families implementing $\mathrm{ADD}^n$, which takes
two integers in the standard $n$-bit representation and returns their sum in the
standard $(n+1)$-bit representation: the ripple-carry adder ($|C_n| = 5n$) and the
Kogge--Stone parallel-prefix adder ($|C_n| = 9, 18, 27, 39, \dots$), for
$n = 2,\dots,9$, uniformly random inputs, and gate-by-gate execution [Eq.~(26)]. Both
families use the same basis with priors: every gate of a given type has the prior given
by that type's output distribution on independent uniform inputs,
$q_{\mathrm{AND}}(1) = 1/4$, $q_{\mathrm{OR}}(1) = 3/4$, $q_{\mathrm{XOR}}(1) = 1/2$.
Both families scale linearly with circuit size but with different slopes: about
0.97 bits per gate for the ripple-carry adder and 0.84 bits per gate for the
Kogge--Stone adder, whose deep prefix levels contain many nearly deterministic gates. At
equal input size, the ripple-carry adder is cheaper (e.g.\ 43.5 versus 77.0 bits at
$n = 9$) because it has far fewer gates. With uniform priors, every gate would
contribute exactly one bit and both families would lie on the line
$\mathcal{MC}(C_n) = |C_n|$. Mismatch cost in bits.}

% Option B: priors matched to each circuit (ripple_vs_kogg2_matched.pdf)
\caption{Mismatch cost of two circuit families implementing $\mathrm{ADD}^n$, which takes
two integers in the standard $n$-bit representation and returns their sum in the
standard $(n+1)$-bit representation: the ripple-carry adder ($|C_n| = 5n$) and the
Kogge--Stone parallel-prefix adder ($|C_n| = 9, 18, 27, 39, \dots$), for
$n = 2,\dots,9$, uniformly random inputs, and gate-by-gate execution [Eq.~(26)]. Each
circuit uses the priors that minimize its own mismatch cost for this input
distribution, $q_{g|\mathrm{pa}(g)} = p_g$, so that the total mismatch cost equals
$\sum_g S(X_g)$. The mismatch cost per gate rises with $n$ for the ripple-carry adder
(from 0.79 to 0.89 bits) and falls for the Kogge--Stone adder (from 0.88 to 0.71 bits),
so the two families separate when plotted against circuit size. The Kogge--Stone curve
is concave because its group-propagate signals over $k$ bits equal 1 with probability
$2^{-k}$: gates in deep prefix levels are nearly deterministic and contribute little. At
equal input size, the ripple-carry adder is cheaper (e.g.\ 40.1 versus 63.7 bits at
$n = 9$) because it has far fewer gates. With uniform priors, every gate would
contribute exactly one bit and both families would lie on the line
$\mathcal{MC}(C_n) = |C_n|$. Mismatch cost in bits.}

% Optional companion (ripple_vs_kogg2_vs_n.pdf or ripple_vs_kogg2_matched_vs_n.pdf)
\caption{The data of Fig.~7 plotted against the input size $n$ (bits per operand). The
ripple-carry adder has the lower mismatch cost for every $n \ge 3$.}

% ----------------------------------------------------------------------
% fig08: Homogeneous vs heterogeneous priors (AND tree with one OR per level)
\caption{Mismatch cost scaling for a circuit family with homogeneous (orange) versus
heterogeneous (blue) priors, for uniformly random inputs and gate-by-gate execution
[Eq.~(26)]. The circuit family consists of binary trees of AND and OR gates in which each
level has one OR gate and the remaining gates are AND gates, so the number of OR gates
equals the depth, $\log_2(|C_n| + 1)$. With identical uniform priors for all gates
(orange), every gate contributes exactly one bit and the mismatch cost equals the circuit
size, $\mathcal{MC}(C_n) = |C_n|$. When the OR gates instead have prior
$q_{\mathrm{OR}}(1) = 0.01$ while the AND gates keep the uniform prior (blue), each OR
gate contributes about 4 additional bits, giving 5.0, 11.4, 19.9 and 32.3 bits for
1, 3, 7 and 15 gates. The mismatch cost then grows as a weighted gate count
(Theorem~2) rather than linearly with circuit size. Mismatch cost in bits.}

% ----------------------------------------------------------------------
% fig09: AND binary trees under agreeing, uniform and disagreeing priors
\caption{Mismatch cost for a family of binary-tree AND circuits under different gate
priors, for a uniform input distribution and gate-by-gate execution [Eq.~(26)].
Because the inputs are re-sampled on every run, the value a gate holds just before it
updates is its output from the previous run, which is independent of its new inputs
[Eq.~(21)]; the mismatch cost of the update measures how well the gate's prior matches
this pre-update state. The priors are: (blue) a prior that agrees with the pre-update
state, $q_{g|\mathrm{pa}(g)} = p_g$, for which the total mismatch cost reduces to
$\sum_g S(X_g)$; (orange) a uniform prior, for which every gate contributes exactly one
bit; (green) a prior that assigns probability 0.97 to the output opposite to that of an
ideal AND gate given its inputs (`disagrees'). The agreeing prior gives the lowest
mismatch cost (7.9 bits for 15 gates, compared with 15.0 bits for the uniform prior and
58.4 bits for the disagreeing prior). Mismatch cost in bits.}

% ----------------------------------------------------------------------
% fig10: Mixed binary-tree family (CONST0, AND, NOR, NAND) under an AND-favouring prior
\caption{Mismatch cost of a mixed binary tree circuit family. This figure illustrates that
the linear cost--size relation need not hold for circuit families in the finite input regime.
The circuit family is constructed as follows: starting with a depth-1 tree which consists of
a CONST gate that always outputs zero, the depth-2 tree is made of AND gates, depth-3 tree is
made from NOR gates and lastly, depth 4 tree is a tree of NAND gates. All gates are evaluated
under the same prior: inputs are assumed independent and uniformly distributed, so that
$q_{\mathrm{pa}(g)}(\boldsymbol{x}_{\mathrm{pa}(g)}) = \tfrac{1}{4}$ for all
$\boldsymbol{x}_{\mathrm{pa}(g)} \in \{0,1\}^2$, and
$q_{g|\mathrm{pa}(g)}(\boldsymbol{x}_g \mid \boldsymbol{x}_{\mathrm{pa}(g)}) =
\begin{cases} 1-\varepsilon & \text{if } \boldsymbol{x}_g = x \land y, \\ \varepsilon & \text{otherwise}, \end{cases}$
with $\varepsilon = 10^{-12}$. This means that the prior strongly favors AND behavior, and
deviations from AND incur mismatch cost. Jointly, $q_{\mathrm{pa}(g)}$ and $q_{g|\mathrm{pa}(g)}$
determine the prior
$q_{g,\mathrm{pa}(g)}(\boldsymbol{x}_g, \boldsymbol{x}_{\mathrm{pa}(g)}) =
q_{g|\mathrm{pa}(g)}(\boldsymbol{x}_g \mid \boldsymbol{x}_{\mathrm{pa}(g)})\,
q_{\mathrm{pa}(g)}(\boldsymbol{x}_{\mathrm{pa}(g)})$.}

% ----------------------------------------------------------------------
% fig11: Single gates: mismatch cost vs gate prior
\caption{Mismatch cost of individual gates as a function of their prior distributions. The
input distribution to the gate is uniform. The x-axis is the prior probability
$q_g(x_g = 1)$. Each curve is minimized when the prior equals the gate's output distribution
($q_g(x_g = 1) = 1/4$ for AND, $3/4$ for OR and $1/2$ for XOR).}

% ----------------------------------------------------------------------
% fig12: Four-layer AND/OR circuit: total MMC vs OR-gate prior
\caption{Total mismatch cost of a circuit with 4 layers, each containing two AND gates
and one OR gate, as a function of the prior of the OR gates, for uniformly random inputs
and gate-by-gate execution [Eq.~(26)]. The circuit has three inputs, and each layer
computes $\mathrm{AND}(n_0, n_1)$, $\mathrm{AND}(n_1, n_2)$ and $\mathrm{OR}(n_2, n_0)$ of
the three nodes $(n_0, n_1, n_2)$ of the previous layer. The prior of the AND gates is
fixed and uniform, while the prior of the OR gates, $q_g(x_g = 1)$, is varied from 0.01
to 0.99. At $q_g(x_g = 1) = 0.5$ all priors are uniform and the total mismatch cost
equals the circuit size, 12 bits. The minimum, 11.2 bits, occurs at
$q_g(x_g = 1) = 0.75$, the probability with which every OR gate in this circuit outputs 1;
the cost rises steeply as the OR prior moves away from it (27.9 bits at 0.01 and
14.7 bits at 0.99). Mismatch cost in bits.}
